\subsection{Средняя скорость}
%%%%%%%%%%%%%%%%%%%%%%%%%%%%%%%%%%%%%%%%%%%%%%%%%%%%
%%%%%%%%%%%%%%%%%%%%%%%%%%%%%%%%%%%%%%%%%%%%%%%%%%%%
%%%%%%%%%%%%%%%%%%%%%%%%%%%%%%%%%%%%%%%%%%%%%%%%%%%%
\z{Катер половину пути плыл со скоростью $v$, а вторую половину — со скоростью $2v$, и его средняя скорость оказалась равной $48\,\text{км/ч}$. Чему равнялась скорость катера на первой половине пути?}
%
{$v = 36\,\text{км/ч}$}

%%%%%%%%%%%%%%%%%%%%%%%%%%%%%%%%%%%%%%%%%%%%%%%%%%%%
\z{Автомобиль проехал половину пути со скоростью $v_1 = 60\,\text{км/ч}$, оставшуюся часть пути он половину времени ехал со скоростью $v_2 = 20\,\text{км/ч}$, а последний участок — со скоростью $v_3 = 100\,\text{км/ч}$. Найдите среднюю скорость автомобиля на всём пути.}
%
{$60\,\text{км/ч}$}

%%%%%%%%%%%%%%%%%%%%%%%%%%%%%%%%%%%%%%%%%%%%%%%%%%%%
\z{Треть всего пути автомобиль проехал с постоянной скоростью $v_1 = 40\,\text{км/ч}$. Затем треть всего времени он ехал с постоянной скоростью $v_2 = 90\,\text{км/ч}$. Найдите среднюю скорость на всем пути, если она оказалась равна средней скорости на оставшемся участке.}
%
{$v_\text{ср}=\sqrt{v_1v_2}=60 \,\text{км/ч}$}

%%%%%%%%%%%%%%%%%%%%%%%%%%%%%%%%%%%%%%%%%%%%%%%%%%%%
\z{Автобус выходит из пункта $A$ и проходит расстояние $s = 40\,\text{км}$ до пункта $B$ со средней скоростью $v_1 = 40\,\text{км/ч}$ и останавливается там на время $t = 20\,\text{мин}$. Затем он возвращается в пункт $A$, проходя расстояние $s$ со средней скоростью $v_2 = 60\,\text{км/ч}$. Найдите среднюю $v_{\text{ср}}$ и среднепутевую $v_{\text{сп}}$ скорости за всё время движения автобуса.}
%
{$v_{\text{ср}} = 0,\quad v_{\text{сп}} = 40\,\text{км/ч}$}

%%%%%%%%%%%%%%%%%%%%%%%%%%%%%%%%%%%%%%%%%%%%%%%%%%%%
\z{Первую треть времени точка движется со скоростью $v_1$, вторую треть — со скоростью $v_2$, последнюю треть — со скоростью $v_3$. Найдите среднюю скорость точки за всё время движения.}
%
{$v_{\text{ср}} = \frac{v_1 + v_2 + v_3}{3}$}

%%%%%%%%%%%%%%%%%%%%%%%%%%%%%%%%%%%%%%%%%%%%%%%%%%%%
\z{Первую треть пути точка движется со скоростью $v_1$, вторую треть — со скоростью $v_2$, последнюю треть пути — со скоростью $v_3$. Найдите среднюю скорость точки за весь путь.}
%
{$v_{\text{ср}} = \frac{3}{\frac{1}{v_1} + \frac{1}{v_2} + \frac{1}{v_3}}$}

%%%%%%%%%%%%%%%%%%%%%%%%%%%%%%%%%%%%%%%%%%%%%%%%%%%%
\z{Турист первую треть всего времени движения шел по лесу на юг со скоростью $v_1 = 3\,\text{км/ч}$, затем треть всего пути перемещался по полю на восток со скоростью $v_2$ и, наконец, по кратчайшему пути по просеке вернулся в исходную точку. Вычислите среднюю путевую скорость $v_0$ туриста. Укажите минимальное возможное значение скорости $v_2$.}
%
{$v_0=\frac{4}{3}v_1=4\,\text{км/ч};$~$v_2=\frac{2}{3}v_1=2\,\text{км/ч}$}

